% generated by roam.py handbook export -- не редактировать, править заметки
% 20261010170941-идея_поиска_оптимального_шага_для_неточнои_оптимизации.org (body 27c455fda576)

\begin{concept}[{Идея поиска оптимального шага для неточной оптимизации}]\label[concept]{note:BE21881B-6DC7-45CB-94F5-9F2E94C0584F}
  В некоторых случаях, когда мы доказываем сходимость какого-то вида \hyperref[note:BEEE81AD-13FC-4846-A160-C72E31760BEE]{град спуска},
  мы получаем оценку
  \[
  f(x_{k+1}) \le f(x_k) - c(\alpha) \cdot r_k^2
  \]

  Где $r_k$~--- \hyperref[note:C42D2C19-B35F-4C0E-AFEB-D880CBB72176]{невязка}, $\alpha$~--- шаг, а $c(\alpha)$~--- некая константа, которая связана с шагом.
  По сути мы выводим \emph{оценку сверху} для следующего значения функции ($f(x_{k+1})$),
  которая связана с шагом градиентного спуска.

  А далее,  мы хотим эту оценку сверху сделать как можно лучше. То есть, уменьшить ее значение.
  Минимизация оценки сверху сводится к максимизации $c(\alpha)$, так как он стоит перед отрицательным слагаемым.

  По такой эвристике мы выбираем оптимальный шаг на каждой итерации.

  \emph{TODO} Обобщить эту идею дальше
\end{concept}

\begin{remark}
  Это можно назвать \hyperref[note:A818191F-B1A6-492C-B38D-5369D169161C]{неточной оптимизацией}, так как вместо значения функции мы минимизируем ее потолок на каждой итерации.
\end{remark}

\paragraph*{Почему мы так делаем?}
Улучшая оценку сверху для шага по направлению, мы решаем задачу оптимизации $c(\alpha)$, а не $f(x_k + \alpha d_k)$.  И мы строим оценку так, чтобы найти $\max_{\alpha} c(\alpha)$ было проще, чем $\min_{\alpha} f(x_k + \alpha d_{k})$ (как в \hyperref[note:8EDDABAC-D2BE-4EC9-8297-DBAC52F70D28]{Точная одномерная оптимизация})

\paragraph*{Где встречается}
\proofpart{Липшицев градиент}
Как пример можно взять \hyperref[note:FFEA1F1D-B70C-4B4A-9386-8344A2A3BBC4]{оценку убывания за шаг градиентного спуска}.
Там $c(\alpha) = \alpha (1- \frac{L\alpha}{2})$~--- это просто парабола, можно найти ее вершину

\proofpart{Липшицев градиент и сильная выпуклость}
Если посмотреть на \hyperref[note:85D9D77D-13A2-42E0-BF93-AD085E0134B9]{Сходимость градиентного спуска для сильно выпуклых функций с липшицевым градиентом}
